\section{Численный метод прямой задачи}
\label{sec:forward-method}

\draft{Применяется метод конечных объёмов~\cite{samarskii2003}.
Производная по времени берётся при движении ячейки с материалом;
теплофизические коэффициенты уточняются итерациями Пикара. Для неявного Эйлера:}
\begin{equation}
 (\rho c)_P^{n+1}|P^{n+1}|
       \frac{T_P^{n+1}-T_P^n}{\Delta t}
       +\sum_{\sigma\subset\partial P}F_{P\sigma}^{n+1}
       =A_P^{n+1}|P^{n+1}|,
 \qquad F_{P\sigma}=-\int_\sigma\lambda\nabla T\cdot\boldsymbol n_P\dd\sigma.
\end{equation}
\draft{На горизонтальных гранях используется двухточечная аппроксимация,
на вертикальных подгранях --- многоточечная. Причина показана на
рис.~\ref{fig:flux}: разность T в центрах с разными глубинами не определяет
горизонтальную производную. Рабочий вариант --- линейная MPFA;
нелинейная NMPFA используется для сопоставления. Многоточечные и монотонные
схемы рассматриваются в~\cite[гл.~2]{inmost2020}; свойства конкретной реализации
подтверждаются отдельно.}
\reportfigure{04_multipoint_flux.png}
  {Две соседние колонки с несовпадающими уровнями центров и разными
  теплопроводностями. Для поля $T=z$ горизонтальный поток должен быть нулевым.}
  {fig:flux}

\draft{Табличный нижний поток задан на единицу горизонтальной проекции:
его вклад для колонки ширины $\Delta x$ равен $q\Delta x$.
Временные шаги попадают в границы эпох. Линейные системы решаются
GMRES(100)~\cite{saad1986} с ILU(2); относительные допуски линейного
и нелинейного решений равны $10^{-8}$. Начальная стационарная задача
использует тот же пространственный оператор.}

\draft{На шаге температура в точке расчёта Ro линейна по времени.
Интегрирование начинается с $b=\max(t^n,t_{\mathrm{dep}})$; $T_a=T(b)$,
$T_b=T(t^{n+1})$, $\beta_i=E_i/R_g$. Первообразная выражается через
интегральную экспоненту; при $b\geq t^{n+1}$ интегралы не изменяются:}
\begin{gather}
 G(\chi)=\frac{e^{-\chi}}{\chi}+\operatorname{Ei}(-\chi),\qquad
 \Phi_i(T)=\beta_iG(\beta_i/T),\\
 \Delta I_i=\frac{A_i(t^{n+1}-b)}{T_b-T_a}
             \bigl(\Phi_i(T_b)-\Phi_i(T_a)\bigr),\quad T_a\ne T_b,\\
 \Delta I_i=A_i(t^{n+1}-b)e^{-\beta_i/T_a},\quad T_a=T_b.
\end{gather}
\draft{Для ускорения табулируется $H(\chi)=\chi^2e^\chi G(\chi)$:
8192 равномерных значения на $[20,200]$ с линейной интерполяцией.
Тогда $\Phi_i(T)=T^2e^{-\beta_i/T}H(\beta_i/T)/\beta_i$.
Точность и стоимость этого вычисления проверяются в разделах
\ref{sec:verification} и~\ref{sec:performance}.}
